% `tex.web` §143, §175 and §897: a word hyphenated where it had a ligature.
%
% §897's `reconstitute` puts a word back together around each break it is
% offered, and the two halves of the resulting discretionary are re-ligatured
% material. §143 says what a ligature node holds — the characters it was made
% of — and §175 prints *those*, then skips the `replace` nodes the two halves
% stand in for:
%
%     disc_node: begin short_display(pre_break(p));
%       short_display(post_break(p));
%       n:=replace_count(p);
%       while n>0 do
%         begin if link(p)<>null then p:=link(p); decr(n); end;
%       end;
%
% Two faults, and each hid the other.
%
% The halves were rebuilt without the record, so they printed the ligature's
% own glyph position rather than its letters — for cmr10's `fi` that is
% character 12, which prints as nothing at all. And the `replace` nodes were
% not skipped, so the unbroken word was printed *as well as* the two halves.
% Together `office` broken after `of` read `of-ffices`: the post-break
% invisible, the whole `ffi` ligature printed in its place. One fault alone
% would have read `of-fiffices` or `of-ces`, and the pair looked like a
% plausible word.
%
% A hyphenation exception rather than patterns, so the case says what it means
% and does not depend on a pattern file.
\catcode`\{=1 \catcode`\}=2
\lccode`\o=`\o \lccode`\f=`\f \lccode`\i=`\i
\lccode`\c=`\c \lccode`\e=`\e \lccode`\s=`\s \lccode`\x=`\x
\lefthyphenmin=1 \righthyphenmin=1
\hyphenation{of-fices}
\tracingonline=1
\tracingparagraphs=1
\hsize=40pt \parindent=0pt \pretolerance=-1 \tolerance=9999
\font\rip=cmr10 \rip
\hyphenchar\rip=`\-

xx offices\par

% §896 decides whether a word may be hyphenated at all, from its first letter:
%
%     if lc_code(c)=0 then goto done1;
%     if (lc_code(c)=c)or(uc_hyph>0) then goto done2 else goto done1;
%
% A character with no `\lccode` is not a letter and stops the word; one whose
% `\lccode` is not itself is a capital, and a capitalised word is hyphenated
% only when `\uchyph` is positive. And the patterns are matched against the
% `\lccode` image, not Unicode's own lowercase — `\lccode` is how a document
% says which characters are letters and what each folds to.
%
% We asked neither question, so `\uchyph=0` hyphenated anyway. trip.tex sets it
% to zero at line 200 for exactly this, and to one on the line after to show
% the difference.
\lccode`\O=`\o

\uchyph=1
xx Offices\par

\uchyph=0
xx Offices\par
\end
